%!TEX root = ../thesis.tex

\chapter{Project Development: Ideal-MHD Solver}

\section{Development priorities after Report 1}
\label{sec:ch2-priorities}

Report~1 supplied three design decisions beyond its Euler baseline. Its
saved-output audit found no strict-IEEE CPU/GPU difference for matched HLLC
cases, so Report~2 retained matched configurations and ported HLL to CUDA. Its
variation matrix found fast-math perturbations and a separate HLLC--Rusanov
method effect, motivating recorded build semantics and an HLL/HLLD method
axis. Finally, checkpoint drift and reference comparisons showed respectively
that final states can miss temporal behaviour and that fp32--fp64 discrepancy
is not reference error. Report~2 therefore used independently executed,
predeclared stopping times and kept discrepancy and error distinct.

The project selected changes by balancing representative physics, numerical
risk, hardware comparability, runtime and evidence value. Brio--Wu supplied a
one-dimensional MHD Riemann problem; Orszag--Tang and Kelvin--Helmholtz
exercised two-dimensional wave interaction and shear evolution. HLL provided a
common CPU/CUDA implementation axis, while CPU HLLD isolated the additional
wave structure and branching of the more elaborate solver. Chapter~3 records
the resulting comparison matrix.

\section{Ideal-MHD and GLM additions}
\label{sec:ch2-mhd-glm}

The MHD extension retained the finite-volume MUSCL construction of
\citet{vanLeer1979muscl} and the Hancock predictor. In nondimensional units
with magnetic permeability set to one, the ideal-MHD equations augmented by
the generalised Lagrange multiplier (GLM) scalar $\psi$ are
\begin{align}
 \partial_t\rho+\nabla\!\cdot(\rho\mathbf v)&=0, \notag\\
 \partial_t(\rho\mathbf v)+\nabla\!\cdot
 [\rho\mathbf v\mathbf v+(p+\tfrac12|\mathbf B|^2)\mathbf I
   -\mathbf B\mathbf B]&=0, \notag\\
 \partial_t\mathbf B+\nabla\!\cdot(\mathbf v\mathbf B-\mathbf B\mathbf v)
   +\nabla\psi&=0, \qquad \nabla\!\cdot\mathbf B=0, \notag\\
 \partial_tE+\nabla\!\cdot[(E+p+\tfrac12|\mathbf B|^2)\mathbf v
   -(\mathbf v\!\cdot\!\mathbf B)\mathbf B]&=0, \notag\\
 \partial_t\psi+c_h^2\nabla\!\cdot\mathbf B&=-\psi/\tau.
 \label{eq:ch2-ideal-mhd}
\end{align}
The ideal-gas closure is $p=(\gamma-1)[E-\rho|\mathbf v|^2/2-
|\mathbf B|^2/2]$. The conserved GLM state is
\begin{equation}
 \mathbf{U}=(\rho,\rho\mathbf{v},\mathbf{B},E,\psi)^{T},\qquad
 E=\frac{p}{\gamma-1}+\frac{1}{2}\rho|\mathbf{v}|^{2}
   +\frac{1}{2}|\mathbf{B}|^{2}.
 \label{eq:ch2-mhd-state}
\end{equation}
The implemented $E$ excludes $\psi$ energy; $\mathbf v$ and $\mathbf B$ retain
three components on one- and two-dimensional grids.

The two-dimensional solver reuses the normal flux by rotating states into the
$y$-normal frame. The GLM flux uses $F(B_n)=\psi$ and
$F(\psi)=c_h^2B_n$. Damping is applied once per full Lie-split $x$--$y$ step:
\begin{equation}
 \psi^{n+1}=\begin{cases}
 \psi^{*}\exp(-\Delta t\,c_h/c_r),&c_h>0\ \text{and}\ c_r>0,\\
 \psi^{*},&c_h\leq0\ \text{or}\ c_r\leq0.
 \end{cases}
 \label{eq:ch2-glm-damping}
\end{equation}
Here $c_h$ is the stepwise maximum directional fast-wave signal speed, while
$c_r$ sets $c_p^2=c_hc_r$ and, for positive parameters, $\tau=c_r/c_h$.
The two-dimensional default $c_r=0.18$ is the ratio $c_p^2/c_h$ identified by
\citet{dedner2002glm} as a compromise between damping and stability. The guard
in Equation~\ref{eq:ch2-glm-damping} makes $c_r=0$ an explicit undamped control:
$\psi$ is left unchanged and no division is evaluated. GLM transports and
damps discrete divergence error but does not enforce a solenoidal field.
The fixed $x$-then-$y$ Lie composition is first-order as a multidimensional
split. Minmod slopes reconstruct the nine conservative components. If a
reconstructed face becomes non-physical, both faces revert to the cell average;
if a line update remains non-physical, that line is recomputed with first-order
HLL fluxes. A post-step failure of the finite, positive-density and
positive-pressure checks terminates the run with a diagnosed numerical failure
and non-zero exit code. No density or pressure floor is applied.

Eight-wave source formulations transport constraint error with the flow
\citep{powellEtAl1999adaptiveMhd}, whereas constrained-transport methods evolve
face-centred magnetic fluxes to preserve a discrete divergence
\citep{evansHawley1988ct,balsaraSpicer1999ct,londrilloDelZanna2004uct,
gardinerStone2005ct}; constrained hyperbolic cleaning refines the GLM source
itself \citep{triccoPrice2012divb}. GLM was retained because it fits the cell-centred,
directionally split code.

Figure~\ref{fig:ch2-implementation-delta} summarises these additions.

\begin{figure}[htbp]
  \centering
  \begin{tikzpicture}[
      box/.style={draw=fpDarkBlue, rounded corners, align=center,
                  minimum height=0.75cm, text width=3.15cm, font=\footnotesize},
      arrow/.style={->, thick, draw=fpGray}]
    \node[box, text width=4.2cm] (base) at (0,0)
      {Report~1 Euler harness\\configuration, reconstruction and output};
    \node[box] (state) at (-3.8,-2.15)
      {Nine-variable ideal-MHD state\\MHD flux and fast waves};
    \node[box] (glm) at (0,-2.15)
      {Rotated two-dimensional sweeps\\GLM transport and damping};
    \node[box] (flux) at (3.8,-2.15)
      {HLL baseline\\CPU HLLD analysis path};
    \node[box] (device) at (-2.0,-4.55)
      {CPU execution and OpenMP build control\\tested CUDA HLL mirror};
    \node[box] (gates) at (2.0,-4.55)
      {Physical-state, completion\\and regression checks};
    \draw[arrow] (base) -- (state);
    \draw[arrow] (base) -- (glm);
    \draw[arrow] (base) -- (flux);
    \draw[arrow] (state) -- (device);
    \draw[arrow] (glm) -- (device);
    \draw[arrow] (glm) -- (gates);
    \draw[arrow] (flux) -- (gates);
  \end{tikzpicture}
  \caption[Ideal-MHD implementation delta]{\textbf{Ideal-MHD implementation
  delta.} Added numerical, execution and acceptance components.}
  \label{fig:ch2-implementation-delta}
\end{figure}

\section{HLL and HLLD solver paths}
\label{sec:ch2-riemann}

HLL uses the left and right MUSCL--Hancock predicted face states. For
$k\in\{L,R\}$, its fast magnetosonic speed and Davis-type bounds are
\begin{align}
 a_k^2&=\gamma p_k/\rho_k,&
 c_{A,k}^2&=|\mathbf B_k|^2/\rho_k,\notag\\
 c_{An,k}^2&=B_{n,k}^2/\rho_k,\notag\\
 c_{f,k}^2&=\tfrac12\!\left[a_k^2+c_{A,k}^2+
 \sqrt{(a_k^2+c_{A,k}^2)^2-4a_k^2c_{An,k}^2}\right],\notag\\
 S_L^{H}&=\min(v_{n,L}-c_{f,L},v_{n,R}-c_{f,R},-c_h),&
 S_R^{H}&=\max(v_{n,L}+c_{f,L},v_{n,R}+c_{f,R},c_h).
 \label{eq:ch2-hll-speeds}
\end{align}
Thus the fan contains the GLM waves $\pm c_h$. Where local fast speeds fall
well below that stepwise maximum, the clamp widens the fan beyond a local
estimate, adding dissipation. With $\mathbf F_k=\mathbf F(\mathbf U_k)$, the
standard flux is
\begin{equation}
 \mathbf F_{H}=\begin{cases}
 \mathbf F_L,&0\le S_L^{H},\\
 \dfrac{S_R^{H}\mathbf F_L-S_L^{H}\mathbf F_R+
 S_L^{H}S_R^{H}(\mathbf U_R-\mathbf U_L)}{S_R^{H}-S_L^{H}},
   &S_L^{H}<0<S_R^{H},\\
 \mathbf F_R,&S_R^{H}\le0.
 \end{cases}
 \label{eq:ch2-hll-flux}
\end{equation}
This compact two-wave solver \citep{hartenLaxVanLeer1983} is the common
CPU/CUDA MHD path. Whenever $c_h>0$ the clamp forces
$S_L^{H}\le-c_h<0<c_h\le S_R^{H}$, so the outer branches of
Equation~\ref{eq:ch2-hll-flux} are unreachable and strict inequalities there
change nothing. HLLD has no such clamp, so its $S_L^*$, $S_R^*$ and $S_M$
ties remain reachable.

HLLD, or HLL with Discontinuities, instead uses the plain Davis outer speeds
\begin{equation}
 S_L=\min(v_{n,L}-c_{f,L},v_{n,R}-c_{f,R}),\qquad
 S_R=\max(v_{n,L}+c_{f,L},v_{n,R}+c_{f,R}),
 \label{eq:ch2-hlld-outer-speeds}
\end{equation}
without the $\pm c_h$ clamp. Following \citet{miyoshiKusano2005hlld}, its
five-wave fan $S_L$, $S_L^*$, $S_M$, $S_R^*$, $S_R$ restores the Alfv\'en and
contact discontinuities, yielding the four intermediate fluxes
$\mathbf F_L^*$, $\mathbf F_L^{**}$, $\mathbf F_R^{**}$, $\mathbf F_R^*$.
The GLM pair is solved first:
\begin{equation}
 B_n^*=\tfrac12(B_{n,L}+B_{n,R})-\frac{\psi_R-\psi_L}{2c_h},\qquad
 \psi^*=\tfrac12(\psi_L+\psi_R)-\frac{c_h}{2}(B_{n,R}-B_{n,L}).
 \label{eq:ch2-hlld-glm-star}
\end{equation}
Every returned flux re-imposes $F(B_n)=\psi^*$ and
$F(\psi)=c_h^2B_n^*$. For $p_T=p+|\mathbf B|^2/2$,
$m_k=S_k-v_{n,k}$ and $D=m_R\rho_R-m_L\rho_L$, the key intermediate speeds are
\begin{align}
 S_M&=\frac{m_R\rho_Rv_{n,R}-m_L\rho_Lv_{n,L}-p_{T,R}+p_{T,L}}{D},&
 \rho_k^*&=\frac{\rho_km_k}{S_k-S_M},\notag\\
 S_L^*&=S_M-\frac{|B_n^*|}{\sqrt{\rho_L^*}},&
 S_R^*&=S_M+\frac{|B_n^*|}{\sqrt{\rho_R^*}}.
 \label{eq:ch2-hlld-speeds}
\end{align}
Fourteen guards test non-finite data, intermediate states and fluxes,
non-positive $\rho_k^*$, and $D$, $S_k-S_M$ and the Alfv\'en denominators
against a relative $64\varepsilon_{\mathrm{mach}}$ floor. Here
$\varepsilon_{\mathrm{mach}}$ is machine epsilon for the active precision; the
multiplier scales the local denominator terms, with a minimum scale of one.
Failure sends only
that interface
to HLL on the same GLM-split states. Otherwise, matched HLL/HLLD runs share
reconstruction, updates and cleaning; the selected Riemann flux is the method
axis, with local fallback retained as part of HLLD.

\section{CPU, OpenMP build control, and CUDA implementation}
\label{sec:ch2-execution}

OpenMP is a shared-memory threading interface, but the MHD CPU path has no
work-sharing directives, so requests of 1, 2, 4 or 8 threads cannot alter the
decomposition. Reduction order could not change the result even under work
sharing: the only reduction here, the stepwise maximum for the CFL condition
and $c_h$, is exactly associative on finite IEEE-754 values. Order dependence
would need an accumulating reduction, absent outside diagnostics. These runs
are a null control.

CUDA, NVIDIA's GPU platform, provides an independent path mirroring CPU HLL
semantics in fp32 and fp64: boundaries, conservative minmod/Hancock
reconstruction, GLM cleaning, sweep and step order, the
Courant--Friedrichs--Lewy (CFL) time-step calculation and per-cell reversion.
CUDA obtains $c_h$ by a two-stage parallel maximum tree rather than a
sequential scan, which associativity again leaves harmless.

The CUDA MHD kernel translation unit is compiled with \texttt{--fmad=false}
whenever CUDA is enabled, scoped to that file, because MSVC's default
/fp:precise does not contract multiply--add on x64
\citep{microsoftMsvcFp2021}. The other relevant nvcc defaults
(\texttt{--ftz=false}, \texttt{--prec-div=true}, \texttt{--prec-sqrt=true})
already conform to IEEE-754 \citep{whiteheadFitFlorea2011gpuFp}, leaving
contraction the only device
transformation to suppress. Agreement is therefore a property of the declared
build configuration rather than of the hardware, and the option is exposed so
that restoring the nvcc default can be measured
(Section~\ref{sec:ch5-hardware}) instead of assumed.

CUDA omits HLLD, per-line fallback and post-step rejection.
\citet{bardDorelli2014gpu} supplies context for a related GPU scheme, not a
performance baseline across different hardware and optimisation choices.

\section{Testing and acceptance criteria}
\label{sec:ch2-gates}

Unit tests isolate state conversion, pressure, wave-speed estimates, fluxes
and directional symmetry. Integration tests exercise complete CPU and CUDA
updates, including GLM cleaning, boundaries, transverse invariance and
supported two-dimensional cases. Regression tests preserve configuration and
output contracts. This division separates faults in numerical operations from
failures of a complete update.

A run enters measurement only if it reaches the requested final time with a
positive step count, finite states, positive density and pressure, fresh output
and completion metadata, so incomplete or non-physical runs are rejected before
aggregation. Property checks such as transverse invariance assess mathematical
behaviour, whereas bitwise CPU/GPU identity assesses consistency across
hardware; neither supplies a numerical or physical reference solution.
